% GHR Calculus for Octonions: Approach and Findings
% Octonion Computation Substrate — Phase 2
%
% Compile with: pdflatex ghr-calculus-octonions.tex (run twice for citations/TOC)
% Requires: amsmath, amssymb, amsthm, mathtools, booktabs, hyperref, geometry

\documentclass[11pt]{article}

\usepackage[margin=1.2in]{geometry}
\usepackage{amsmath, amssymb, amsthm}
\usepackage{mathtools}
\usepackage{booktabs}
\usepackage{hyperref}
\usepackage{microtype}
\usepackage{parskip}

% Theorem environments
\newtheorem{theorem}{Theorem}[section]
\newtheorem{proposition}[theorem]{Proposition}
\newtheorem{lemma}[theorem]{Lemma}
\newtheorem{definition}[theorem]{Definition}
\newtheorem{remark}[theorem]{Remark}

% Notation
\newcommand{\OO}{\mathbb{O}}
\newcommand{\HH}{\mathbb{H}}
\newcommand{\CC}{\mathbb{C}}
\newcommand{\RR}{\mathbb{R}}
\newcommand{\ee}[1]{\mathbf{e}_{#1}}
\newcommand{\conj}[1]{\overline{#1}}
\newcommand{\norm}[1]{\lVert #1 \rVert}
\newcommand{\abs}[1]{\lvert #1 \rvert}
\newcommand{\ip}[2]{\langle #1,\, #2 \rangle}
\newcommand{\Jac}[2]{\mathbf{J}_{#1}[#2]}
\newcommand{\Lmat}[1]{L_{#1}}
\newcommand{\Rmat}[1]{R_{#1}}
\newcommand{\C}[3]{C_{#1#2}^{#3}}
\newcommand{\pd}[2]{\frac{\partial #1}{\partial #2}}
\newcommand{\defeq}{\coloneqq}

\title{GHR Calculus Extended to Octonions:\\
Approach, Implementation, and Findings}
\author{Octonion Computation Substrate}
\date{March 2026}

\begin{document}
\maketitle
\tableofcontents

\bigskip

% ============================================================
\section{Introduction}
% ============================================================

Gradient-based optimisation of octonionic neural networks requires a rigorous
theory of differentiation in the octonionic algebra $\OO$.  For quaternionic
networks, the \emph{Generalized Hamilton--Real} (GHR) calculus of Xu \&
Mandic~\cite{xu2015ghr} provides a principled framework: Wirtinger-like
derivative pairs, product and chain rules, and gradient expressions that
preserve algebraic structure.  Extending GHR to $\OO$ is non-trivial because
octonion multiplication is \emph{non-associative}, and the GHR framework for
$\HH$ relies on associativity in a fundamental way.

This document describes the approach taken, the key theoretical obstacle
(the rotation ambiguity arising from non-associativity), its resolution via
the Moufang identity, the resulting implementation of analytic Jacobians and
parenthesization-aware chain rule, and the quantitative verification results.

% ============================================================
\section{Background: GHR Calculus for Quaternions}
% ============================================================

\subsection{Wirtinger Derivatives in \texorpdfstring{$\CC$}{C}}

For a function $f:\CC\to\CC$ written in real coordinates as
$f(x,y)=u+iv$, the Wirtinger (Cauchy--Riemann) derivative pair is
\begin{equation}
  \pd{f}{z} = \tfrac{1}{2}\!\left(\pd{f}{x} - i\,\pd{f}{y}\right),
  \qquad
  \pd{f}{\conj{z}} = \tfrac{1}{2}\!\left(\pd{f}{x} + i\,\pd{f}{y}\right).
\end{equation}
The complex analyticity condition $\partial f/\partial\conj{z}=0$ recovers the
Cauchy--Riemann equations.  For a real-valued loss $L$, the gradient direction
for descent is $\partial L/\partial\conj{z}$ (the conjugate derivative).

\subsection{GHR Derivatives for Quaternions}

For a quaternion $\mathbf{q} = q_a + q_b\,\mathbf{i} + q_c\,\mathbf{j} + q_d\,\mathbf{k}$
and a unit quaternion frame $\boldsymbol{\mu}$, the GHR rotation is
$\mathbf{q}_\mu \defeq \boldsymbol{\mu}\,\mathbf{q}\,\boldsymbol{\mu}^{-1}$.
The GHR derivatives are then defined as
\begin{equation}
  \pd{f}{\mathbf{q}_\mu}
  = \tfrac{1}{4}\!\left(
      \pd{f}{q_a}
    - \pd{f}{q_b}\,\mathbf{i}_\mu
    - \pd{f}{q_c}\,\mathbf{j}_\mu
    - \pd{f}{q_d}\,\mathbf{k}_\mu
    \right),
\end{equation}
\begin{equation}
  \pd{f}{\mathbf{q}_\mu^*}
  = \tfrac{1}{4}\!\left(
      \pd{f}{q_a}
    + \pd{f}{q_b}\,\mathbf{i}_\mu
    + \pd{f}{q_c}\,\mathbf{j}_\mu
    + \pd{f}{q_d}\,\mathbf{k}_\mu
    \right),
\end{equation}
where $(\mathbf{i}_\mu, \mathbf{j}_\mu, \mathbf{k}_\mu)$ is the rotated
imaginary basis.  For real-valued loss functions, Proposition~4.8 of
Xu \& Mandic~\cite{xu2015ghr} shows that left and right derivatives coincide,
and the optimization gradient is $\partial L/\partial\mathbf{q}^*_\mu$.

The GHR framework also supplies a product rule: for
$f(\mathbf{q}) = g(\mathbf{q})\,h(\mathbf{q})$,
\begin{equation}
  \pd{f(\mathbf{q})}{\mathbf{q}_\mu}
  = g(\mathbf{q})\,\pd{h(\mathbf{q})}{\mathbf{q}_\mu}
  + \pd{g(\mathbf{q})}{\mathbf{q}_{h(\mathbf{q})\cdot\mu}}\,h(\mathbf{q}),
\end{equation}
where the rotation frame for $g$ is shifted by the value of $h$.  This product
rule exploits quaternion associativity in the rotation $\mathbf{q}_\mu = \boldsymbol{\mu}\,\mathbf{q}\,\boldsymbol{\mu}^{-1}$.

% ============================================================
\section{Extending GHR to Octonions}
% ============================================================

\subsection{The Non-Associativity Obstruction}

The octonions $\OO \cong \RR^8$ have a basis
$\{\ee{0}, \ee{1}, \ldots, \ee{7}\}$ with multiplication encoded by
integer-valued structure constants $C_{ij}^k \in \{-1,0,+1\}$ satisfying
\begin{equation}
  \ee{i}\,\ee{j} = \sum_{k=0}^{7} C_{ij}^k\,\ee{k}.
\end{equation}
The Cayley--Dickson structure constants (Baez 2002, mod-7 Fano plane basis) are
used throughout.  Unlike $\HH$, octonion multiplication is \emph{not}
associative: in general $(a\,b)\,c \neq a\,(b\,c)$.

The GHR rotation $R_\mu(\mathbf{o}) \defeq \boldsymbol{\mu}\,\mathbf{o}\,\boldsymbol{\mu}^{-1}$
requires specifying a parenthesization:
$(\boldsymbol{\mu}\,\mathbf{o})\,\boldsymbol{\mu}^{-1}$
vs.\
$\boldsymbol{\mu}\,(\mathbf{o}\,\boldsymbol{\mu}^{-1})$.
In $\HH$, associativity makes these equal; in $\OO$, they may differ.

\subsection{Resolution via the Moufang Identity}

The octonions, while non-associative, are an \emph{alternative} algebra:
every subalgebra generated by two elements is associative.  Alternative algebras
satisfy the Moufang identities, of which the most useful for our purposes is:
\begin{equation}\label{eq:moufang}
  (\boldsymbol{\mu}\,\mathbf{o})\,\boldsymbol{\mu}^{-1}
  = \boldsymbol{\mu}\,(\mathbf{o}\,\boldsymbol{\mu}^{-1}).
\end{equation}
The key observation is that $\boldsymbol{\mu}$ appears \emph{twice} (as
$\boldsymbol{\mu}$ and $\boldsymbol{\mu}^{-1}$), and the Moufang identity
guarantees that any expression of the form $xyx$ is well-defined regardless of
parenthesization in an alternative algebra.  Therefore, the octonionic rotation
$R_\mu(\mathbf{o}) = \boldsymbol{\mu}\,\mathbf{o}\,\boldsymbol{\mu}^{-1}$ is
\emph{unambiguous} despite non-associativity, and the GHR framework for $\HH$
can be transported to $\OO$ without modification of the rotation definition.

\subsection{The Octonionic GHR Derivative Pair}

For an octonion $\mathbf{o} = o_0 + \sum_{k=1}^{7} o_k\,\ee{k}$, the natural
extension of the GHR derivative pair uses 8 real partial derivatives:
\begin{align}
  \pd{f}{\mathbf{o}}
  &= \frac{1}{8}\!\left(
        \pd{f}{o_0} - \sum_{k=1}^{7}\pd{f}{o_k}\,\ee{k}
     \right),\label{eq:ghr-full}\\[4pt]
  \pd{f}{\mathbf{o}^*}
  &= \frac{1}{8}\!\left(
        \pd{f}{o_0} + \sum_{k=1}^{7}\pd{f}{o_k}\,\ee{k}
     \right).\label{eq:ghr-conj}
\end{align}
The factor $\tfrac{1}{8}$ (compared to $\tfrac{1}{4}$ for $\HH$ and
$\tfrac{1}{2}$ for $\CC$) reflects the 8-dimensional nature of $\OO$ and
ensures the fundamental reconstruction identity
\begin{equation}
  df = \pd{f}{\mathbf{o}}\,d\mathbf{o} + \pd{f}{\mathbf{o}^*}\,d\mathbf{o}^*
\end{equation}
holds under the standard octonionic inner product.

For a real-valued loss $L$ and octonionic parameter $\mathbf{o}$, the
optimization gradient is
\begin{equation}
  \nabla_{\mathbf{o}} L
  = \pd{L}{\mathbf{o}^*}
  = \frac{1}{8}\!\left[
      \pd{L}{o_0},\;
      \pd{L}{o_1},\;
      \ldots,\;
      \pd{L}{o_7}
    \right]^\top,
\end{equation}
which is simply the real gradient vector $\nabla_{\RR^8} L$ scaled by
$\tfrac{1}{8}$.  The GHR formalism adds structural content for
\emph{non-real-valued} intermediate computations by tracking both derivatives
of the pair and supplying algebraically consistent product and chain rules.

% ============================================================
\section{Implementation: Jacobian-First Approach}
% ============================================================

\subsection{Motivation}

A direct symbolic derivation of the GHR product rule for octonions requires
careful book-keeping of rotation frames that shift with intermediate function
values.  We adopt instead a \emph{Jacobian-first} strategy: ground every
primitive operation in an explicit $8 \times 8$ real Jacobian matrix, verify it
against finite differences, and build composites via the standard matrix-product
chain rule.  This sidesteps non-associativity at the single-operation level
(a Jacobian is a linear map on $\RR^8$, and matrix multiplication is
associative), while correctly encoding non-associativity at the composition
level through explicit parenthesization tracking.

\subsection{Structure Constants and Multiplication Jacobian}

The product $(a \cdot b)_k = \sum_{i,j} C_{ij}^k\,a_i\,b_j$ is bilinear, so
its Jacobians are linear in the \emph{other} factor:
\begin{align}
  \Jac{a}{k,i}
  &= \pd{(a\cdot b)_k}{a_i}
   = \sum_j C_{ij}^k\,b_j
   = (\Rmat{b})^\top_{ki},\label{eq:J-a}\\[4pt]
  \Jac{b}{k,j}
  &= \pd{(a\cdot b)_k}{b_j}
   = \sum_i C_{ij}^k\,a_i
   = (\Lmat{a})^\top_{kj},\label{eq:J-b}
\end{align}
where $\Lmat{a}$ and $\Rmat{b}$ are the left and right multiplication matrices
already present in the codebase.  In \texttt{einsum} notation:
\begin{equation}
  \mathbf{J}_a = \texttt{einsum}(\text{``ijk, ...j} \to \text{...ki''}, C, b),\qquad
  \mathbf{J}_b = \texttt{einsum}(\text{``ijk, ...i} \to \text{...kj''}, C, a).
\end{equation}

\subsection{Jacobians for All Seven Primitives}

Analytic $8\times 8$ Jacobians are derived for each primitive operation.

\paragraph{Conjugation.}
$\conj{\mathbf{o}} = (o_0, -o_1, \ldots, -o_7)^\top$, so
$\mathbf{J}_{\text{conj}} = \operatorname{diag}(1,-1,-1,-1,-1,-1,-1,-1)$.

\paragraph{Inverse.}
$\mathbf{o}^{-1} = \conj{\mathbf{o}}/\norm{\mathbf{o}}^2$.  By the quotient rule:
\begin{equation}
  \Jac{}{k,i}
  = \frac{(J_\text{conj})_{ki}}{\norm{\mathbf{o}}^2}
  - \frac{2\,\conj{o}_k\,o_i}{\norm{\mathbf{o}}^4}.
\end{equation}

\paragraph{Exponential.}
Writing $\mathbf{o} = a + \mathbf{v}$ with $a = o_0$, $\mathbf{v} = (o_1,\ldots,o_7)$, and
$r = \norm{\mathbf{v}}$:
\begin{equation}
  \exp(\mathbf{o}) = e^a\!\left(\cos r + \frac{\sin r}{r}\,\mathbf{v}\right).
\end{equation}
The Jacobian has block structure: a $1\times 1$ scalar--scalar block, a
$1\times 7$ real--imaginary block, a $7\times 1$ imaginary--real block, and a
$7\times 7$ imaginary--imaginary block involving a rank-one outer product
$\mathbf{v}\mathbf{v}^\top$ plus a diagonal term:
\begin{equation}
  \mathbf{J}_{\mathrm{im,im}}
  = e^a\!\left[
      \frac{\cos r\cdot r - \sin r}{r^3}\,\mathbf{v}\mathbf{v}^\top
      + \frac{\sin r}{r}\,\mathbf{I}_7
    \right].
\end{equation}
Near $r=0$, L'H\^opital limits apply: $\sin r/r \to 1$ and
$(\cos r \cdot r - \sin r)/r^3 \to -1/3$.

\paragraph{Logarithm.}
With $q = \norm{\mathbf{o}}$ and $\theta = \arccos(a/q)$:
\begin{equation}
  \log(\mathbf{o}) = \log q + \frac{\theta}{r}\,\mathbf{v}.
\end{equation}
The Jacobian has a similar block structure.  The outer-product coefficient
$(a/q^2 - \theta/r)/r^2$ admits a finite limit as $r \to 0$ (derived via Taylor
expansion of $\theta$ about $r = 0$).

\paragraph{Inner product and cross product.}
The inner product $\ip{a}{b} = \sum_i a_i b_i$ has trivial Jacobians
$\mathbf{J}_a = b^\top$ and $\mathbf{J}_b = a^\top$.  The cross product
$\text{cross}(a,b) = \operatorname{Im}(\operatorname{Im}(a) \cdot \operatorname{Im}(b))$
has Jacobians equal to the $7\times 7$ imaginary--imaginary block of the
multiplication Jacobians computed at the purified arguments, with the real row
and column zeroed.

\subsection{Triple Verification}

Every primitive is checked against three independent gradient computations:
\begin{enumerate}
  \item \textbf{Analytic Jacobian} (closed-form derivation above).
  \item \textbf{Numeric Jacobian} (central finite differences with
        $\varepsilon = 10^{-7}$ at \texttt{float64}).
  \item \textbf{Autograd Jacobian} (column-wise backward passes through the
        \texttt{torch.autograd.Function} implementation).
\end{enumerate}
All three agree to relative error below $10^{-5}$ (see Section~\ref{sec:results}).

\subsection{Higher-Order Derivatives}

Each \texttt{torch.autograd.Function} backward pass is implemented
using only differentiable PyTorch primitives (\texttt{einsum},
\texttt{exp}, \texttt{cos}, \texttt{sin}, etc.), and \emph{only inputs} are
saved in \texttt{ctx.save\_for\_backward}.  This makes the backward pass
itself differentiable, enabling \texttt{create\_graph=True} for second-order
gradient computation (Hessian eigenspectrum analysis in Phase~5).

% ============================================================
\section{Parenthesization-Aware Chain Rule}
% ============================================================

\subsection{The Non-Associativity of Compositions}

For standard $\RR^n$ functions, the chain rule gives a Jacobian product
$\mathbf{J}_\text{total} = \mathbf{J}_\text{outer} \cdot \mathbf{J}_\text{inner}$,
and matrix multiplication is associative so the order of application is
irrelevant.  For octonionic compositions, a crucial distinction arises: the
\emph{Jacobians} compose associatively (they are linear maps on $\RR^8$), but
the \emph{operations being differentiated} are not associative.  That is,
$(a \cdot b) \cdot c$ and $a \cdot (b \cdot c)$ are genuinely different functions
with genuinely different Jacobians, even though the composition of their
Jacobians via matrix multiplication is well-defined in either order.

\subsection{Binary Tree Representation}

Each multiplication chain is represented as an explicit binary tree:
\begin{itemize}
  \item A \textbf{Leaf} node holds an operand index.
  \item A \textbf{Node} holds an operation label (\texttt{mul}, \texttt{exp},
        \texttt{log}, \texttt{conjugate}, \texttt{inverse}) and references to
        left and right subtrees.
\end{itemize}
For $n$ operands in a product chain, there are $C_{n-1}$ distinct binary tree
structures, where $C_k$ is the $k$-th Catalan number:
$C_0=1, C_1=1, C_2=2, C_3=5, C_4=14, C_5=42, \ldots$

The function \texttt{all\_parenthesizations(n)} generates all $C_{n-1}$ trees
by the standard recursive construction: split the chain at position $k$
($1 \le k < n$), enumerate all trees for each sub-chain, and combine.

\subsection{Autograd Integration}

Each \texttt{evaluate\_tree} call dispatches to \texttt{OctonionMulFunction.apply}
(or the corresponding unary function) at each internal node.  PyTorch's
autograd engine records the exact computation graph, including which operands
were multiplied in which order.  The backward pass therefore propagates
gradients through the correct tree structure automatically, without any
special-casing in the backward code.

\subsection{The Naive Baseline}

A \emph{naive} chain rule is defined as the one that always parenthesizes
left-to-right: $((x_0 \cdot x_1) \cdot x_2) \cdot \cdots$.  For a real-valued
loss, the gradient of the naive rule is the Jacobian of this specific function.
When the true computation uses a different parenthesization, the naive gradient
is simply \emph{wrong}: it computes the gradient of a different function.

% ============================================================
\section{Octonionic Analyticity}
% ============================================================

\subsection{CR-Like Conditions}

In complex analysis, $f:\CC\to\CC$ is holomorphic iff its $2\times 2$ real
Jacobian equals a left-multiplication matrix $\Lmat{c}$ for some
$c\in\CC$.  The octonionic analogue defines $f:\OO\to\OO$ to be
\emph{octonionic-analytic} (in the Fueter/left-regular sense) iff its
$8\times 8$ real Jacobian satisfies
\begin{equation}
  \mathbf{J} = \Lmat{c},
  \qquad\text{where}\quad
  \Lmat{c}[k,j] = \sum_i c_i\,C_{ij}^k.
\end{equation}
The \emph{CR residual} is $\norm{\mathbf{J} - \Lmat{c}}_F$ where
$c = \mathbf{J}[:,0]$ (first column of $\mathbf{J}$, since $\Lmat{c}\,\ee{0} = c$).

\subsection{Classification of Primitives}

The analyticity conditions are extremely restrictive in $\OO$, far more
so than in $\CC$ or $\HH$.

\begin{center}
\begin{tabular}{lll}
  \toprule
  Operation & Analytic? & Reason \\
  \midrule
  $f(\mathbf{o}) = c \cdot \mathbf{o}$ (left multiplication) & Yes & $\mathbf{J} = \Lmat{c}$ exactly \\
  $f(\mathbf{o}) = \alpha\,\mathbf{o}$, $\alpha\in\RR$ & Yes & $\mathbf{J} = \alpha\,\mathbf{I} = \Lmat{\alpha\ee{0}}$ \\
  $f(\mathbf{o}) = \mathbf{o}$ (identity) & Yes & $\mathbf{J} = \mathbf{I} = \Lmat{\ee{0}}$ \\
  $f(\mathbf{o}) = \mathbf{o} \cdot c$ (right multiplication) & No & $\Rmat{c} \neq \Lmat{c}$ in general \\
  $f(\mathbf{o}) = \conj{\mathbf{o}}$ (conjugation) & No & $\mathbf{J} = \operatorname{diag}(1,-1,\ldots,-1)$ \\
  $\exp$, $\log$ & No & Complex block Jacobian structure \\
  \bottomrule
\end{tabular}
\end{center}

The scarcity of analytic functions is a known feature of octonionic analysis and
contrasts sharply with complex analysis.  Non-associativity severely restricts
which local linear approximations can be represented as multiplication by a fixed
octonion.

% ============================================================
\section{Quantitative Results}
\label{sec:results}
% ============================================================

\subsection{Gradient Accuracy: All 14 Parenthesizations}

We tested all $C_4 = 14$ parenthesizations of 5-operand product chains against
finite-difference Jacobians at \texttt{float64} precision.

\begin{center}
\begin{tabular}{rll}
  \toprule
  Tree index & Tree structure & Max relative error \\
  \midrule
  0  & $(x_0(x_1(x_2(x_3 x_4))))$               & $1.04 \times 10^{-8}$ \\
  1  & $(x_0(x_1((x_2 x_3)x_4)))$                & $2.95 \times 10^{-8}$ \\
  2  & $(x_0((x_1 x_2)(x_3 x_4)))$               & $3.44 \times 10^{-8}$ \\
  3  & $(x_0((x_1(x_2 x_3))x_4))$                & $5.12 \times 10^{-6}$ \\
  4  & $(x_0(((x_1 x_2)x_3)x_4))$                & $2.37 \times 10^{-8}$ \\
  5  & $((x_0 x_1)(x_2(x_3 x_4)))$               & $8.68 \times 10^{-7}$ \\
  6  & $((x_0 x_1)((x_2 x_3)x_4))$               & $1.35 \times 10^{-8}$ \\
  7  & $((x_0(x_1 x_2))(x_3 x_4))$               & $1.47 \times 10^{-7}$ \\
  8  & $(((x_0 x_1)x_2)(x_3 x_4))$               & $1.50 \times 10^{-8}$ \\
  9  & $((x_0(x_1(x_2 x_3)))x_4)$                & $4.95 \times 10^{-7}$ \\
  10 & $((x_0((x_1 x_2)x_3))x_4)$                & $7.35 \times 10^{-8}$ \\
  11 & $(((x_0 x_1)(x_2 x_3))x_4)$               & $1.24 \times 10^{-7}$ \\
  12 & $(((x_0(x_1 x_2))x_3)x_4)$                & $1.26 \times 10^{-8}$ \\
  13 & $((((x_0 x_1)x_2)x_3)x_4)$                & $2.06 \times 10^{-7}$ \\
  \midrule
  \multicolumn{2}{l}{All patterns passed ($<10^{-5}$)} & \textbf{yes} \\
  \bottomrule
\end{tabular}
\end{center}

The worst-case error is $5.12 \times 10^{-6}$ (tree index~3), well within the
$10^{-5}$ relative-error criterion.  Thirteen of the fourteen patterns are
within $10^{-6}$.  The maximum Jacobian norm difference \emph{between} distinct
parenthesizations of the same five operands is $14.63$, confirming that different
parenthesizations define quantifiably different functions.

\subsection{Naive vs.\ Correct Chain Rule}

To quantify the cost of ignoring non-associativity, we compare the parenthesization-aware
Jacobian against the \emph{naive} left-associative chain rule across 1{,}000
random inputs at each depth.

\begin{center}
\begin{tabular}{rrlll}
  \toprule
  Depth & \# Patterns & Mean $\norm{\Delta J}$ & Std & Max $\norm{\Delta J}$ \\
  \midrule
  3 & 2   & $11.14$ & $3.38$ & $21.49$ \\
  5 & 14  & $31.01$ & $14.34$ & $85.69$ \\
  7 & 132 & $73.23$ & $43.09$ & $248.35$ \\
  \bottomrule
\end{tabular}
\end{center}

The mean gradient error grows roughly as depth$^{1.4}$, compounding rapidly.
The mean cosine similarity between naive and correct gradient vectors at depth~3
is $0.331$ — the vectors point in substantially different directions.  A network
trained with the naive chain rule would follow entirely wrong gradient directions
for any computation involving 3 or more chained octonionic multiplications.

\subsection{GPU/CPU Parity}

All 8 backward-pass tests (covering \texttt{mul}, \texttt{exp}, \texttt{log},
\texttt{conjugate}, \texttt{inverse}, \texttt{OctonionLinear}, batched
multiplication, and a 3-operand composition) were verified on an AMD Radeon
RX~7900~XTX (gfx1100, ROCm~7.2) against CPU computation at \texttt{float64}.
Maximum elementwise difference: $< 10^{-12}$ in all cases.

% ============================================================
\section{Tradeoffs and Design Decisions}
% ============================================================

\subsection{Native GHR vs.\ Pseudo-Real Matrix Representation}

Two principled approaches exist for octonionic differentiation:

\begin{enumerate}
  \item \textbf{Native GHR (this work).}  Derive derivatives directly in $\OO$,
        producing octonionic-valued Wirtinger derivatives and a product rule that
        respects algebraic structure.  Requires resolving the rotation ambiguity
        (Section~3.2).
  \item \textbf{Pseudo-real matrix representation.}  Embed each octonion as an
        $8\times 8$ real matrix via left multiplication, reduce the problem to
        real matrix calculus, and apply standard automatic differentiation.  This
        approach (used in e.g.\ Octonion Phase Retrieval, 2023~\cite{opr2023}) is
        straightforward but reduces the derivative to $\RR^8$ --- the algebraic
        structure of $\OO$ is lost, and the result is equivalent to treating the
        network as a real network with structured weight sharing.
\end{enumerate}

The native approach was chosen because (a) the Wirtinger derivative pair is the
mathematical object needed for gradient-based optimization while preserving
algebraic structure, (b) the pseudo-real approach gives no additional insight
into how non-associativity affects gradient flow, and (c) the parenthesization
analysis (a potential research contribution) requires the native representation.

\subsection{Jacobian-First vs.\ Symbolic GHR}

A fully symbolic derivation would apply the GHR product rule
directly to expressions like $f(\mathbf{o}) = (\mathbf{w} \cdot \mathbf{o}) \cdot \mathbf{b}$,
producing closed-form octonionic gradient expressions without going through the
$\RR^8$ Jacobian.  The Jacobian-first approach was chosen instead for three
reasons:

\begin{enumerate}
  \item \textbf{Verifiability.}  An $8\times 8$ matrix can be checked against
        finite differences column by column; a symbolic octonionic expression
        cannot be directly verified without computing the same matrix.
  \item \textbf{Generality.}  The Jacobian-first approach handles arbitrary
        computation graphs; symbolic derivations must be re-done for each
        new formula.
  \item \textbf{Associativity of composition.}  The chain rule for Jacobian
        matrices is $\mathbf{J}_\text{total} = \mathbf{J}_n \cdots \mathbf{J}_1$,
        and matrix multiplication \emph{is} associative even though the
        operations themselves are not.  This cleanly separates the
        non-associativity problem (which parenthesization was used in the
        forward pass) from the derivative computation (which always reduces to
        associative matrix products).
\end{enumerate}

The cost is a slight loss of algebraic elegance: gradients are $\RR^8$ vectors
reinterpreted as octonions, rather than octonionic expressions derived from
first principles in $\OO$.

\subsection{Fused vs.\ Primitive-Decomposed Jacobians}

All composite Jacobians are built from the seven primitives via the chain rule,
rather than deriving fused Jacobians (e.g., a single $8\times 8$ matrix for
$\exp(a \cdot b)$).  This is more modular and easier to verify independently,
but may introduce small numerical errors relative to a fused derivation because
intermediate Jacobians are computed at intermediate points rather than directly
at the composed function's arguments.  In practice the errors are negligible
at \texttt{float64} (see Section~\ref{sec:results}).

\subsection{Number of Wirtinger Derivatives}

Complex numbers require 2 Wirtinger derivatives
($\partial f/\partial z$, $\partial f/\partial\conj{z}$); quaternions require
4 in the full GHR framework (one pair per rotation direction $\mu$, though for
real-valued loss the pairs coincide).  For octonions with the Moufang-resolved
rotation, the natural extension gives a single pair
$(\partial f/\partial\mathbf{o},\, \partial f/\partial\mathbf{o}^*)$,
corresponding to the identity rotation frame $\mu = \ee{0}$.

In principle, 7 additional rotation frames (one per imaginary unit) could
each generate an independent Wirtinger pair, yielding 16 derivative components.
However, for optimization with a real-valued loss, Proposition~4.8 of
Xu \& Mandic~\cite{xu2015ghr} generalizes directly: left and right derivatives
coincide, and the gradient direction is fully determined by the conjugate
derivative $\partial L/\partial\mathbf{o}^*$, which is the $\RR^8$ gradient
scaled by $\tfrac{1}{8}$.  The additional rotation-dependent pairs encode
the structure of $f$ as an octonionic function but are not needed for
gradient descent.

% ============================================================
\section{API Summary}
% ============================================================

The implementation is in \texttt{src/octonion/calculus/} and exports the
following public interface:

\begin{center}
\begin{tabular}{ll}
  \toprule
  Symbol & Description \\
  \midrule
  \texttt{ghr\_derivative(partials)} & GHR full derivative \eqref{eq:ghr-full} \\
  \texttt{conjugate\_derivative(partials)} & GHR conjugate derivative \eqref{eq:ghr-conj} \\
  \texttt{wirtinger\_from\_jacobian(J)} & Convert $8\times 8$ Jacobian to Wirtinger pair \\
  \texttt{jacobian\_mul(a, b)} & Analytic Jacobians for $a \cdot b$ \\
  \texttt{jacobian\_exp(o)}, \texttt{jacobian\_log(o)} & Analytic Jacobians for exp, log \\
  \texttt{jacobian\_conjugate(o)}, \texttt{jacobian\_inverse(o)} & Analytic Jacobians for conj, inv \\
  \texttt{jacobian\_inner\_product(a,b)}, \texttt{jacobian\_cross\_product(a,b)} & Analytic Jacobians \\
  \texttt{numeric\_jacobian(fn, x)} & Central finite-difference Jacobian \\
  \texttt{octonion\_gradcheck(fn, inputs)} & Custom gradient checker \\
  \texttt{Leaf}, \texttt{Node}, \texttt{all\_parenthesizations(n)} & Binary tree types and enumeration \\
  \texttt{evaluate\_tree(tree, operands)} & Autograd-tracked tree evaluation \\
  \texttt{CompositionBuilder} & High-level composition API \\
  \texttt{is\_octonionic\_analytic(fn, x)} & CR residual test \\
  \texttt{cauchy\_riemann\_octonion(J)} & CR residual from Jacobian \\
  \texttt{suggest\_lr(model, loader)} & GHR gradient magnitude heuristic \\
  \bottomrule
\end{tabular}
\end{center}

% ============================================================
\section{Conclusion}
% ============================================================

We have successfully extended the GHR calculus to the octonions.  The central
theoretical obstacle --- the ambiguity of the rotation $R_\mu(\mathbf{o}) =
\boldsymbol{\mu}\,\mathbf{o}\,\boldsymbol{\mu}^{-1}$ under non-associativity
--- is resolved by the Moufang identity, which guarantees that expressions
with a repeated element are parenthesization-independent in alternative algebras.

The implementation uses a Jacobian-first strategy: closed-form $8\times 8$
real Jacobians for each primitive operation, verified triple-redundantly against
finite differences and autograd.  Non-associativity is handled explicitly through
binary tree representations of computation graphs; the parenthesization-aware
chain rule produces correct gradients for all $C_{n-1}$ parenthesizations of
$n$-operand chains.

Quantitatively, the naive left-associative chain rule produces gradients with
a mean angular error of $\cos^{-1}(0.33) \approx 71^\circ$ from the correct
gradient at depth~3, growing to roughly $10\times$ larger gradient norm errors
at depth~7.  No practical octonionic network could be trained correctly without
parenthesization-aware differentiation.

% ============================================================
\begin{thebibliography}{9}
% ============================================================

\bibitem{xu2015ghr}
  N.~Xu and D.~P.~Mandic,
  ``The theory of quaternion matrix derivatives,''
  \emph{IEEE Trans.\ Signal Process.}, vol.~63, no.~6, pp.~1543--1556, 2015.
  (GHR calculus; product rule, chain rule, Proposition~4.8.)

\bibitem{xu2016ghr2}
  N.~Xu, C.~Cheong, and D.~P.~Mandic,
  ``A new proof and interpretation of the GHR calculus,''
  \emph{Royal Society Open Science}, vol.~3, 2016.

\bibitem{baez2002}
  J.~C.~Baez,
  ``The octonions,''
  \emph{Bull.\ Amer.\ Math.\ Soc.}, vol.~39, pp.~145--205, 2002.
  (Structure constants, Fano plane, Moufang identities.)

\bibitem{parcollet2019}
  T.~Parcollet, M.~Morchid, and G.~Linar\`es,
  ``A survey of quaternion neural networks,''
  \emph{Artif.\ Intell.\ Rev.}, vol.~53, pp.~2957--2982, 2020.

\bibitem{gaudet2018}
  C.~J.~Gaudet and A.~S.~Maida,
  ``Deep quaternion networks,''
  in \emph{Proc.\ IJCNN}, 2018.

\bibitem{opr2023}
  \emph{Octonion Phase Retrieval},
  arXiv:2308.15784, 2023.
  (Pseudo-real matrix approach for octonionic derivatives --- approach not
   adopted here, but validates the problem statement.)

\end{thebibliography}

\end{document}
